% =====================================================================
\chapter{Estimación por máxima verosimilitud y suavizado}
\label{tema:3}
% =====================================================================

Los parámetros de un modelo de $n$-gramas se estiman a partir de un corpus. Este tema demuestra que el estimador de máxima verosimilitud es la frecuencia relativa, analiza su principal defecto (las probabilidades nulas) y presenta dos remedios: el suavizado aditivo y la interpolación. Se sigue~\cite[§3.1.2 y §3.6]{jm3} y, para la interpretación bayesiana, \cite[§2.2]{bishop}.

\section{Verosimilitud de un corpus}

Sea un corpus formado por secuencias, cada una delimitada por \fin\ al principio y al final, y sea $N(u,v)$ el número de veces que el token $v$ sigue inmediatamente al token $u$. Se escribe $N(u)=\sum_{v}N(u,v)$ para el número de veces que $u$ aparece como token anterior.

\begin{definicion}[Verosimilitud]
Dada una matriz de transición $P$, la \emph{verosimilitud} del corpus es el producto de las probabilidades de todos sus bigramas,
\[
\mathcal{L}(P)=\prod_{u,v\in\V}P_{uv}^{\,N(u,v)},
\qquad
\ell(P)=\ln\mathcal{L}(P)=\sum_{u,v\in\V}N(u,v)\ln P_{uv},
\]
y $\ell$ se llama \emph{log-verosimilitud}. El \emph{estimador de máxima verosimilitud} (EMV) es la matriz estocástica que maximiza $\ell$.
\end{definicion}

El logaritmo convierte productos en sumas y evita el subdesbordamiento numérico: con $10^{3}$ factores iguales a $0{,}9$ el producto vale aproximadamente $10^{-46}$.

\section{El estimador de máxima verosimilitud}

\begin{teorema}[EMV del modelo de bigramas]\label{teo:emv}
Supóngase $N(u)>0$ para todo $u$. El estimador de máxima verosimilitud de la matriz de transición es la frecuencia relativa
\[
\widehat{P}_{uv}=\frac{N(u,v)}{N(u)}.
\]
\end{teorema}

\begin{proof}
\emph{Paso 1 (separación por filas).} La función $\ell(P)=\sum_{u}\bigl(\sum_{v}N(u,v)\ln P_{uv}\bigr)$ es suma de términos que dependen cada uno de una sola fila, y las restricciones $\sum_{v}P_{uv}=1$ también afectan a una sola fila. Por tanto, basta maximizar cada fila por separado.

\emph{Paso 2 (lagrangiana).} Fijada una fila $u$, se escribe $N_v=N(u,v)$, $p_v=P_{uv}$ y $n=N(u)$. Se maximiza $f(p)=\sum_{v}N_v\ln p_v$ con la restricción $\sum_v p_v=1$. Si algún $N_v>0$ tuviera $p_v=0$, entonces $f=-\infty$; luego en el máximo $p_v>0$ siempre que $N_v>0$. La lagrangiana es
\[
\Lambda(p,\lambda)=\sum_{v}N_v\ln p_v-\lambda\Bigl(\sum_{v}p_v-1\Bigr).
\]

\emph{Paso 3 (condiciones de primer orden).} $\partial\Lambda/\partial p_v=N_v/p_v-\lambda=0$, de donde $p_v=N_v/\lambda$ para los $v$ con $N_v>0$. Los $v$ con $N_v=0$ no contribuyen a $f$, y cualquier masa asignada a ellos se resta a los demás; por tanto, en el óptimo, $p_v=0$ para ellos.

\emph{Paso 4 (multiplicador).} Imponiendo la restricción, $\sum_{v}N_v/\lambda=1$, luego $\lambda=\sum_v N_v=n$ y $p_v=N_v/n$.

\emph{Paso 5 (máximo global).} $f$ es una suma de funciones cóncavas $N_v\ln p_v$ con $N_v\ge 0$, y el conjunto de distribuciones es convexo; un punto crítico de una función cóncava sobre un convexo es un máximo global.
\end{proof}

\begin{ejemplo}\label{ej:conteos}
Para el corpus «ana», «ala», «luna» los bigramas son 13 y la matriz de conteos es la del \cref{tab:conteos}. Dividiendo cada fila por su total se obtienen las probabilidades de la \cref{fig:cadena}.
\end{ejemplo}

\begin{table}[h]
\centering\small
\caption{Conteos $N(u,v)$ del \cref{ej:conteos} (fila: token anterior; columna: token siguiente).}\label{tab:conteos}
\begin{tabular}{@{}c|ccccc|c@{}}
\toprule
$u\backslash v$ & \fin & \texttt{a} & \texttt{l} & \texttt{n} & \texttt{u} & $N(u)$ \\
\midrule
\fin & 0 & 2 & 1 & 0 & 0 & 3 \\
\texttt{a} & 3 & 0 & 1 & 1 & 0 & 5 \\
\texttt{l} & 0 & 1 & 0 & 0 & 1 & 2 \\
\texttt{n} & 0 & 2 & 0 & 0 & 0 & 2 \\
\texttt{u} & 0 & 0 & 0 & 1 & 0 & 1 \\
\bottomrule
\end{tabular}
\end{table}

\section{Ceros y suavizado aditivo}

El EMV asigna probabilidad nula a todo bigrama que no aparece en el corpus. Si un texto de evaluación contiene uno de ellos, su probabilidad es~0 y la pérdida del \cref{tema:5} es infinita. Este fenómeno es una forma de sobreajuste~\cite[§3.5]{jm3}.

\begin{definicion}[Suavizado aditivo]\label{def:laplace}
Sea $\alpha>0$. El estimador con \emph{suavizado aditivo} es
\[
\widehat{P}_{\alpha}(v\mid u)=\frac{N(u,v)+\alpha}{N(u)+\alpha K}.
\]
Con $\alpha=1$ recibe el nombre de \emph{suavizado de Laplace}; con $0<\alpha<1$, de \emph{suavizado add-$k$}~\cite[§3.6.1–3.6.2]{jm3}.
\end{definicion}

\begin{proposicion}\label{prop:laplace}
Para cada $u$, $\widehat{P}_{\alpha}(\cdot\mid u)$ es una distribución de probabilidad con todas sus entradas estrictamente positivas. Además, coincide con la media a posteriori de la fila $u$ cuando se adopta una distribución a priori de Dirichlet de parámetros $(\alpha,\dots,\alpha)$.
\end{proposicion}

\begin{proof}
Positividad: el numerador es al menos $\alpha>0$. Suma: $\sum_{v}(N(u,v)+\alpha)=N(u)+\alpha K$, que es el denominador. Para la segunda afirmación, la distribución de Dirichlet es conjugada de la multinomial: si la fila $u$ tiene a priori $\mathrm{Dir}(\alpha,\dots,\alpha)$ y se observan los conteos $N(u,\cdot)$, la distribución a posteriori es $\mathrm{Dir}(\alpha+N(u,v_1),\dots,\alpha+N(u,v_K))$, cuya media en la coordenada $v$ es $(N(u,v)+\alpha)/(N(u)+\alpha K)$~\cite[§2.2.1]{bishop}.
\end{proof}

\begin{ejemplo}
En el \cref{ej:conteos} con $\alpha=1$ y $K=5$: $\widehat{P}_1(\texttt{n}\mid\fin)=(0+1)/(3+5)=1/8$, mientras que el EMV daba~0; y $\widehat{P}_1(\texttt{a}\mid\texttt{n})=(2+1)/(2+5)=3/7\approx 0{,}4286$, frente al valor~1 del EMV. El suavizado transfiere masa de los sucesos observados a los no observados.
\end{ejemplo}

\section{Interpolación}

Otra estrategia consiste en combinar estimadores de distinto orden~\cite[§3.6.3]{jm3}.

\begin{definicion}[Interpolación lineal]
Dados un estimador de bigramas $\widehat{P}_2(v\mid u)$ y uno de unigramas $\widehat{P}_1(v)$, y pesos $\lambda_1,\lambda_2\ge 0$ con $\lambda_1+\lambda_2=1$, el estimador \emph{interpolado} es
\[
\widehat{P}_{\mathrm{int}}(v\mid u)=\lambda_2\,\widehat{P}_2(v\mid u)+\lambda_1\,\widehat{P}_1(v).
\]
\end{definicion}

Es una distribución, por ser combinación convexa de distribuciones. Los pesos se eligen maximizando la verosimilitud de un conjunto de validación separado del de entrenamiento.

\begin{lecturas}
\item Jurafsky y Martin~\cite{jm3}, §3.1.2 «How to estimate probabilities», §3.5 «Generalizing vs.\ overfitting the training set» y §3.6 «Smoothing, Interpolation, and Backoff».
\item Bishop~\cite{bishop}, §2.2 «Multinomial variables» y §2.2.1 «The Dirichlet distribution».
\item Chen y Goodman~\cite{chen-goodman}, estudio comparativo de técnicas de suavizado.
\end{lecturas}

\section*{Ejercicios del tema 3}
\addcontentsline{toc}{section}{Ejercicios del tema 3}

\begin{ejercicio}{2}{eje:3-laplace}
Aplique el suavizado de Laplace ($\alpha=1$) al modelo del \cref{eje:2-sol-matriz}. Calcule $P(\fin\mid\fin)$ y $P(\texttt{l}\mid\texttt{a})$ y compárelas con las del EMV.
\end{ejercicio}

\begin{ejercicio}{2}{eje:3-u1}
Corpus de cuatro frases con marcadores $\langle s\rangle$ y $\langle/s\rangle$, que cuentan como tokens del vocabulario: «$\langle s\rangle$ yo soy Ana $\langle/s\rangle$», «$\langle s\rangle$ Ana yo soy $\langle/s\rangle$», «$\langle s\rangle$ yo soy Ana $\langle/s\rangle$» y «$\langle s\rangle$ yo no quiero sopa ni Ana $\langle/s\rangle$». Calcule $P(\text{Ana}\mid\text{soy})$ con un modelo de bigramas suavizado con Laplace.
\fuente{adaptado de Jurafsky y Martin~\cite{jm3}, cap.~3, ejercicio~3.4.}
\end{ejercicio}

\begin{ejercicio}{2}{eje:3-interp}
Con el corpus del \cref{eje:3-u1}, calcule $P(\text{Ana}\mid\text{soy})$ interpolando el bigrama y el unigrama de máxima verosimilitud con $\lambda_1=\lambda_2=1/2$. Para el unigrama, cuente todos los tokens, incluidos $\langle s\rangle$ y $\langle/s\rangle$.
\fuente{adaptado de Jurafsky y Martin~\cite{jm3}, cap.~3, ejercicio~3.7.}
\end{ejercicio}

\begin{ejercicio}{3}{eje:3-limite}
Demuestre que, para $u$ fijo con $N(u)>0$, $\lim_{\alpha\to 0^{+}}\widehat{P}_{\alpha}(v\mid u)$ es el EMV y $\lim_{\alpha\to\infty}\widehat{P}_{\alpha}(v\mid u)=1/K$. Interprete $\alpha$ como un número de «observaciones ficticias».
\end{ejercicio}
